\chapter{Iron Studies}
\section*{What Is This Test?}
Iron studies (also called the iron panel) consist of serum iron, total iron-binding capacity (TIBC), transferrin saturation (serum iron/TIBC $\times$ 100), and ferritin. Together they assess total body iron stores, iron supply for erythropoiesis, and differentiate the major causes of microcytic anemia. Serum iron measures the amount of iron bound to transferrin in the circulation. TIBC measures the total capacity of transferrin to bind iron (a proxy for transferrin concentration). Transferrin saturation is the percentage of transferrin iron-binding sites occupied. Ferritin is the intracellular iron storage protein; serum ferritin concentration directly reflects total body iron stores (1 mcg/L of serum ferritin corresponds to approximately 8--10 mg of stored iron). In iron deficiency anemia (IDA): decreased serum iron, increased TIBC (transferrin synthesis upregulated by iron-responsive element-binding protein/IRE-BP), decreased transferrin saturation (<16\%), decreased ferritin (<15--30 ng/mL) \cite{camaschella2015iron}. In anemia of chronic disease (ACD): decreased serum iron, normal or decreased TIBC, normal or mildly decreased transferrin saturation (15--25\%), increased ferritin (acute phase reactant) \cite{ganz2019anemia}. In hereditary hemochromatosis: increased serum iron, decreased TIBC (or normal), increased transferrin saturation (>45\% in women and >50\% in men is a screening threshold), markedly increased ferritin (>300 ng/mL in men, >200 ng/mL in women).
\section*{Alternative Names}
\begin{itemize}
  \item Iron Panel
  \item Serum Iron and TIBC
  \item Transferrin Saturation
  \item Ferritin
  \item Fe/TIBC
\end{itemize}
\section*{Principle}
Serum iron is measured by colorimetric assay using ferrozine or similar chromogen after release of iron from transferrin by acidification. TIBC is measured by saturating transferrin with excess iron, removing unbound iron, and measuring iron content. Transferrin saturation is calculated. Ferritin is measured by two-site sandwich chemiluminescent immunoassay.
\section*{History}
In 1934, George Minot, William Murphy, and George Whipple received the Nobel Prize in Physiology or Medicine for their work on liver therapy for anemia, which established the central role of iron (and later vitamin B12) in erythropoiesis. Serum iron was first measured colorimetrically in the 1930s--1950s, and transferrin saturation became a standard measure of iron supply to the marrow in the 1960s. In 1972, Addison and colleagues developed the first immunoradiometric assay for serum ferritin, showing that ferritin closely reflects total body iron stores and providing a sensitive blood test for iron deficiency \cite{addison1972immunoradiometric}. In 1996, Feder and colleagues identified the HFE gene, linking C282Y homozygosity to hereditary hemochromatosis and establishing transferrin saturation and ferritin as the standard screening tests. Subsequent guidelines from the American Association for the Study of Liver Diseases and the World Health Organization codified the current diagnostic thresholds for iron deficiency, anemia of chronic disease, and hemochromatosis.
\section*{Why This Test Is Ordered}
\begin{itemize}
  \item Evaluation of microcytic anemia to distinguish iron deficiency anemia from anemia of chronic disease and thalassemia trait
  \item Confirmation of iron deficiency before iron therapy, especially when the CBC and peripheral smear are inconclusive
  \item Investigation of unexplained normocytic or macrocytic anemia with coexistent iron deficiency
  \item Screening for hereditary hemochromatosis in patients with elevated transferrin saturation, elevated ferritin, or suggestive symptoms (fatigue, arthralgia, hepatomegaly), and in first-degree relatives of affected patients \cite{bacon2011diagnosis}
  \item Monitoring of iron stores in chronic kidney disease, in patients on erythropoiesis-stimulating agents, and after bariatric surgery
  \item Monitoring of transfusional iron overload (ferritin $>$1,000 ng/mL) in multiply transfused patients with thalassemia or sickle cell disease
\end{itemize}
\section*{Primary vs Secondary Test}
Iron studies are a primary diagnostic test for iron deficiency and iron overload. Serum ferritin is the single best initial test for iron deficiency \cite{guyatt1992laboratory}, while the combination of transferrin saturation and ferritin is the standard first-line screen for hereditary hemochromatosis. When results are borderline, secondary tests such as soluble transferrin receptor, reticulocyte hemoglobin content, or HFE genotyping are used.
\section*{How This Test Is Performed}
A fasting morning blood sample is collected by standard venipuncture into a serum or lithium-heparin tube. Serum iron is measured by a colorimetric assay (ferrozine or ferene), TIBC by saturating transferrin with excess iron and measuring the bound iron, and ferritin by a two-site sandwich chemiluminescent immunoassay. Transferrin saturation is calculated as serum iron/TIBC $\times$ 100. Results are typically available within 1--2 days \cite{tietz2023textbook}.
\section*{What Preparation Is Needed}
A fasting morning sample is preferred because serum iron has marked diurnal variation (highest in the morning) and falls after iron-rich meals and iron supplements. Patients should avoid iron supplements for at least 12--24 hours before the blood draw. No other preparation is required.
\section*{Contraindications and Precautions}
\begin{itemize}
  \item There are no absolute contraindications. Interpretive cautions:
  \item (1) ferritin is an acute-phase reactant and is falsely elevated by infection, inflammation, and liver disease;
  \item (2) recent blood transfusion or oral iron ingestion falsely elevates serum iron and ferritin;
  \item (3) the soluble transferrin receptor assay may be needed when ferritin falls in the 15--100 ng/mL ``gray zone'';
  \item (4) iron studies are unreliable within 1--2 weeks after acute blood loss because stores are depleted before serum iron falls.
\end{itemize}
\section*{Risks}
Minimal risk. Slight pain or bruising at the venipuncture site. Rarely: hematoma, infection, or vasovagal syncope.
\section*{What Happens After the Test}
Results are interpreted as a panel in the context of the CBC, reticulocyte count, and clinical history. Confirmed iron deficiency is treated with oral or intravenous iron, and ferritin is rechecked in 2--3 months to document repletion. Patients with elevated transferrin saturation and ferritin are referred for HFE genotyping and phlebotomy if hemochromatosis is confirmed.
\section*{Understanding Results}
\subsection*{Below Normal}
Iron deficiency anemia: Serum iron decreased, TIBC increased ($>$400 mcg/dL), transferrin saturation <16\%, ferritin <15--30 ng/mL (ferritin <15 ng/mL is essentially diagnostic of iron deficiency). Anemia of chronic disease: Serum iron decreased, TIBC normal or decreased, ferritin increased (>100--200 ng/mL). Functional iron deficiency in CKD on erythropoietin: ferritin 100--500 ng/mL with transferrin saturation <20\%.
\subsection*{Above Normal}
Hereditary hemochromatosis: Serum iron elevated, TIBC normal or decreased, transferrin saturation >45\% (women) or >50\% (men), ferritin >200/300 ng/mL (screening thresholds). HFE gene mutations (C282Y homozygosity). Transfusional siderosis: ferritin >1,000 ng/mL in multiply transfused patients. Acute liver injury: release of stored iron. Ferritin as acute phase reactant: infection, inflammation, malignancy (ferritin >1,000--10,000 ng/mL).
\section*{Normal Results}
Serum iron: 60--170 mcg/dL (males), 50--150 mcg/dL (females). TIBC: 250--450 mcg/dL. Transferrin saturation: 20--50\% (males), 15--50\% (females). Ferritin: 20--250 ng/mL (females), 30--300 ng/mL (males). Decreased ferritin (<15 ng/mL) is highly specific for iron deficiency. Serum iron has diurnal variation (highest in the morning, lowest in the evening); fasting morning sample preferred. Recent iron ingestion (supplements) will falsely elevate serum iron.
\section*{Follow-up Tests}
Ferritin (most sensitive and specific for iron deficiency), soluble transferrin receptor (sTfR) --- elevated in iron deficiency, normal in ACD (useful for differentiating IDA + ACD), HFE gene testing (C282Y, H63D for hemochromatosis), complete blood count with RBC indices, hemoglobin electrophoresis (beta-thalassemia trait), bone marrow biopsy with Prussian blue iron stain (gold standard for iron stores).
\section*{References}
\bibliographystyle{unsrtnat}
\bibliography{references}
\clearpage
\section*{Regulatory Status and Authorization Date}
This chapter describes a test category, not a specific commercial product. FDA, CDSCO, EU IVDR, and MHRA approval or registration dates therefore cannot be assigned without the assay manufacturer, product name, instrument, and intended use. For laboratory-developed tests, an individual agency approval date may not exist. See the Preface for the interpretation of regulatory status.

\clearpage
